By Theorem \claimref{singularity.rate-discharge-hypothesis-iii}, under (P, \claimref{singularity.p-univalence}) and (R, \claimref{singularity.r-regularity-fixed-locus}) every boundary mode is thermal at $T = \kappa/2\pi$, and $\kappa = 1/4M$ (Proposition \claimref{singularity.plane-rate}), so $T(M) = 1/8\pi M$ and $1/T = 8\pi M$. The Clausius relation along the family, the hypothesis of Hawking's argument, gives $dS/dM = 8\pi M$, and with $S(0) = 0$, $S(M) = 4\pi M^{2}$. The horizon is at $r = 2M$, of area $A = 4\pi(2M)^{2} = 16\pi M^{2}$, so $S = A/4$. Along the family $dA/dM = 32\pi M$ and $\kappa/8\pi = 1/32\pi M$, so $dM = (\kappa/8\pi)\,dA$ identically and $dS = dM/T = (2\pi/\kappa)(\kappa/8\pi)\,dA = dA/4$. If the period were $p$ instead of $2\pi/\kappa$, then $1/T = p$ and $dS = p\,dM = p(\kappa/8\pi)\,dA$, so $S = (p\kappa/2\pi)\,A/4$: the quarter is the period of (R) and nothing else.
